\begin{proposition}[The zero of the Weierstrass function]\hypertarget{tag-01Q9}{}\label{fam:src:B.1.3}\label{fam:st:B.1.3}
\StatementRecord{fam:src:B.1.3}{proved}{01Q9}{\SeriesRef{fam:src:B.1.1}}

For the original affine logarithm $\mu$, we have $\wp_\tau(\mu)=0$ identically. Equivalently the generator $P=[\mu]$ of \ref{fam:src:B.1.1} is one of the two sections $P_*=(0,\sqrt2t^2(t-1))$ and $-P_*$.
\end{proposition}


\begin{proof}
\begin{enumerate}[label=\arabic*.,leftmargin=*]
\item We give the intersection calculation that identifies the generator. At $t=1$, put $u=t-1$ and blow up the origin. In the $u$-chart $x=uX,y=uY$, the strict equation is
$Y^2=uX^3-3(1+u)^3X+2(1+u)^4$. Its exceptional curve is smooth, and $P_*$ meets it at $(X,Y)=(0,\sqrt2)$. In the other chart $u=xU,y=xY$, the equation is $Y^2=x-3(1+xU)^3U+2(1+xU)^4U^2$. The implicit function theorem solves for $x$ at the common point; the two fibre branches $U=0,x=0$ meet with intersection multiplicity two. These charts cover the tangent conic and resolve the singularity. The two components have multiplicity one; the one met by $P_*$ is the nonidentity component $D$, with $D^2=-2$ from intersection with the full fibre.

\item At $t=0$, two blow-ups in the successive $t$-charts give
\[
 Y^2=tX^3+3t^2(1-t)X+2t^2(1-t)^2,
\quad
 Y_2^2=t^2X_2^3+3t(1-t)X_2+2(1-t)^2.
\]
The second equation has two smooth multiplicity-one components $E_{2,\pm}$ at $Y_2=\pm\sqrt2$, and $P_*$ meets one of them. The remaining singular point, in the second $X$-chart, has
\[
 V^2=H(X,T),\quad H=X^2T+(2+3X)T^2-(4X+3X^2)T^3+2X^2T^4,
 \qquad t=XT.
\]
Let $T_0(X)$ solve $H_T(X,T_0)=0$. It has $T_0=-X^2/4+O(X^3)$. Taylor's formula gives
$H-H(X,T_0)=(T-T_0)^2h_0$, $h_0(0,0)=2$, and $H(X,T_0)=X^4e(X)$, $e(0)=-1/8$. Substitution in the displayed polynomial proves these identities with all higher terms retained. With $Z=(T-T_0)\sqrt{h_0}$, $a=V-Z$, $b=(V+Z)/e(X)$, the equation becomes exactly $ab=X^4$ through a local analytic isomorphism. Its resolution has charts
\[
 a=U_i^{i+1}V_i^i,\quad b=U_i^{3-i}V_i^{4-i},\quad X=U_iV_i
 \quad(0\le i\le3),
\]
with changes $U_{i+1}=1/V_i,V_{i+1}=U_iV_i^2$. They give a chain of three $(-2)$-curves. These are the charts of the two successive point blow-ups resolving $ab=X^4$, so the resolution is proper. Retaining $t=XT$ in the comparison gives multiplicities $2,3,2$ on that chain. The two new multiplicity-one components meet its ends; the first exceptional curve meets its middle component, and the original strict fibre meets the other end of that first exceptional curve. All these statements follow by substituting the displayed monomials and their orders; the only fibre components of multiplicity one are the strict component and $E_{2,\pm}$.

\item Order the six nonidentity components by the graph
$\Theta_1-\Theta_2-\Theta_3-\Theta_4$ and $\Theta_2-\Theta_5-\Theta_6$, with $E_{2,\pm}=\Theta_4,\Theta_6$. Their intersection matrix is $-C$, where
\[
 C=\begin{pmatrix}2&-1&0&0&0&0\\-1&2&-1&0&-1&0\\0&-1&2&-1&0&0\\0&0&-1&2&0&0\\0&-1&0&0&2&-1\\0&0&0&0&-1&2\end{pmatrix},\quad
 C^{-1}=\frac13\begin{pmatrix}6&9&6&3&6&3\\9&18&12&6&12&6\\6&12&10&5&8&4\\3&6&5&4&4&2\\6&12&8&4&10&5\\3&6&4&2&5&4\end{pmatrix}.
\]
Multiplication verifies this inverse. Self-intersections follow as well from the complete multiplicity vector and $F\cdot\Theta_i=0$.

\item The form $dx\wedge dt/y$ is holomorphic and nonvanishing on the finite resolved charts. In the last charts its expression is a unit times $dX\wedge da/a=-dU_i\wedge dV_i$; the earlier blow-up equations cancel the apparent denominators at $y=0$. At infinity it is $-v^{-1}dX\wedge dv/Y$ and has one fibre pole. Thus $K_S=-F$. Adjunction gives $Q^2=-1$ for every section $Q$, since $F\cdot Q=1$. A section meets a unique multiplicity-one component transversely at a smooth fibre point: the pullback of the full base equation to the section has order one, excluding all other possibilities.

\item Let $v_Q$ be the intersection vector with $\Theta_i$ and let $d_Q=Q\cdot D$. The possibilities are $v_Q=0,e_4,e_6$ and $d_Q=0,1$. Define the rational divisor class
\[
 \varphi(Q)=Q-O-(Q\cdot O+1)F+\sum\Theta_i(C^{-1}v_Q)_i+\tfrac12d_QD.
\]
It is orthogonal to $O,F,\Theta_i,D$ by direct multiplication. The chord-and-tangent rational function on the generic cubic has divisor $(Q)+(R)-(Q+R)-(O)$; extending it to the surface adds only vertical components. Applying the displayed orthogonal projection therefore proves $\varphi(Q+R)=\varphi(Q)+\varphi(R)$. Expansion gives
\[
 h(Q)=-\varphi(Q)^2=2+2Q\cdot O-v_Q^tC^{-1}v_Q-\tfrac12d_Q.
\]
For $Q\ne O$, the nonnegative intersection $Q\cdot O$ gives $h(Q)\ge2-4/3-1/2=1/6$. The visible $P_*$ is disjoint from $O$ everywhere, and the explicit blow-up charts place it on both nonidentity components. Hence $h(P_*)=1/6$. The logarithm argument in \ref{fam:src:B.1.1} proves $P_*=nP$ for an integer $n\ne0$. Additivity gives $1/6=n^2h(P)\ge n^2/6$, so $n=\pm1$. The pointed cubic map in \ref{fam:src:B.1.1} has $x(P)=f_{\rm sc}^2\wp_\tau(\mu)$ with $f_{\rm sc}\ne0$ on the ordinary locus. This proves the asserted zero identity.
\end{enumerate}
\end{proof}